% preamble-exam-zh.tex — 共享 preamble
% 用于所有 exam-zh 模板
%
% 样式策略：
%   屏幕 → 语义彩色（蓝=关键想法，红=易错，绿=路线，灰=提示/教师备注）
%   打印 → 自动降级为黑白（通过 print mode 或 monochrome 选项）

\documentclass{exam-zh}

% ---------- 基础配置 ----------
\examsetup{
  page/size = a4paper,
  page/show-head = false,
  page/show-foot = false,
  sealline/show = false,
  font = times,
  math-font = xits,
}

\usepackage{tcolorbox}
\usepackage{needspace}
\tcbuselibrary{skins,breakable}

% ---------- 打印/屏幕颜色切换 ----------
% 用法：\documentclass[monochrome]{exam-zh} 即可切换为纯黑白
% 注意：不要使用 \if@xxx 放在 makeatother 外部；这里改成普通条件名。
\newif\ifeduprintcolor
\eduprintcolortrue

\makeatletter
\@ifclasswith{exam-zh}{monochrome}{\eduprintcolorfalse}{}
\makeatother

% ---------- 语义颜色定义 ----------
% 屏幕：有色彩区分；打印：降级为灰度
\ifeduprintcolor
  % --- 屏幕彩色模式 ---
  \definecolor{edu-blue}{RGB}{47, 82, 122}       % 关键想法
  \definecolor{edu-blue-bg}{RGB}{243, 246, 252}
  \definecolor{edu-red}{RGB}{180, 40, 40}         % 易错提醒
  \definecolor{edu-red-bg}{RGB}{255, 245, 240}
  \definecolor{edu-green}{RGB}{34, 100, 60}       % 路线图
  \definecolor{edu-green-bg}{RGB}{242, 250, 245}
  \definecolor{edu-orange}{RGB}{160, 90, 0}       % 训练目标
  \definecolor{edu-orange-bg}{RGB}{255, 248, 235}
  \definecolor{edu-gray}{RGB}{100, 100, 100}      % 提示/教师备注
  \definecolor{edu-gray-bg}{RGB}{248, 248, 248}
  \definecolor{edu-purple}{RGB}{90, 50, 120}      % 思考题
  \definecolor{edu-purple-bg}{RGB}{248, 244, 252}
\else
  % --- 打印黑白模式 ---
  \definecolor{edu-blue}{gray}{0.15}
  \definecolor{edu-blue-bg}{gray}{0.96}
  \definecolor{edu-red}{gray}{0.25}
  \definecolor{edu-red-bg}{gray}{0.97}
  \definecolor{edu-green}{gray}{0.20}
  \definecolor{edu-green-bg}{gray}{0.97}
  \definecolor{edu-orange}{gray}{0.25}
  \definecolor{edu-orange-bg}{gray}{0.98}
  \definecolor{edu-gray}{gray}{0.40}
  \definecolor{edu-gray-bg}{gray}{0.97}
  \definecolor{edu-purple}{gray}{0.30}
  \definecolor{edu-purple-bg}{gray}{0.98}
\fi

% ---------- 自定义教学环境 ----------
% 共性：enhanced + breakable（跨页不断裂）

% 原题卡片 — 强边框，突出显示
\newtcolorbox{problemcard}{
  enhanced, breakable,
  colback=edu-blue-bg,
  colframe=edu-blue,
  boxrule=1.2pt,
  arc=0pt,
  left=3mm, right=3mm, top=3mm, bottom=3mm,
}

% 解题路线图 — 左边线 + 浅底
\newtcolorbox{routemap}{
  enhanced, breakable,
  colback=edu-green-bg,
  colframe=edu-green!30,
  boxrule=0pt,
  borderline west={2.5pt}{0pt}{edu-green},
  left=3mm, right=3mm, top=3mm, bottom=3mm,
}

% 关键想法 — 蓝色左边线
\newtcolorbox{keyideabox}{
  enhanced, breakable,
  colback=edu-blue-bg,
  colframe=edu-blue!30,
  boxrule=0pt,
  borderline west={2.5pt}{0pt}{edu-blue},
  left=3mm, right=3mm, top=3mm, bottom=3mm,
}

% 易错提醒 — 红色左边线
\newtcolorbox{mistakebox}{
  enhanced, breakable,
  colback=edu-red-bg,
  colframe=edu-red!20,
  boxrule=0pt,
  borderline west={3pt}{0pt}{edu-red},
  left=3mm, right=3mm, top=2.5mm, bottom=2.5mm,
}

% 提示 — 灰色虚线左边线
\newtcolorbox{hintbox}{
  enhanced, breakable,
  colback=edu-gray-bg,
  colframe=edu-gray!20,
  boxrule=0pt,
  borderline west={2pt}{0pt}{edu-gray, dashed},
  left=3mm, right=3mm, top=2mm, bottom=2mm,
}

% 思考题 — 紫色虚线左边线
\newtcolorbox{thinkbox}{
  enhanced, breakable,
  colback=edu-purple-bg,
  colframe=edu-purple!20,
  boxrule=0pt,
  borderline west={2pt}{0pt}{edu-purple, dashed},
  left=2.5mm, right=2.5mm, top=2.5mm, bottom=2.5mm,
}

% 教师备注 — 灰色左边线，小字
\newtcolorbox{teachernote}{
  enhanced, breakable,
  colback=edu-gray-bg,
  colframe=edu-gray!15,
  boxrule=0pt,
  borderline west={2pt}{0pt}{edu-gray!60},
  left=2.5mm, right=2.5mm, top=2.5mm, bottom=2.5mm,
  fontupper=\small,
  colupper=black!60,
}

% 训练目标 — 橙色左边线，小字
\newtcolorbox{traininggoal}{
  enhanced, breakable,
  colback=edu-orange-bg,
  colframe=edu-orange!15,
  boxrule=0pt,
  borderline west={1.5pt}{0pt}{edu-orange},
  left=2mm, right=2mm, top=1mm, bottom=1mm,
  fontupper=\small,
}

% 答题区 — 无边框，纯留白
\newcommand{\answerarea}[1][40mm]{%
  \vspace{2mm}%
  \begin{tcolorbox}[
    enhanced, breakable,
    colback=white,
    colframe=white,
    boxrule=0pt,
    height=#1,
    valign=top,
    bottom=0mm,
    after skip=0mm,
  ]
  \end{tcolorbox}%
}

% ---------- 字体设置 ----------
% exam-zh (ctexbook) already sets CJK fonts via ctex.
% Only override if specific fonts are needed.
% macOS: Songti SC, STHeiti, STFangsong
% Linux: SimSun, SimHei, FangSong

% ---------- 页面设置 ----------
% exam-zh (ctexbook) already loads geometry.
% Override margins if needed:
% \geometry{top=16mm, bottom=16mm, left=14mm, right=14mm}

\examsetup{
  question/show-answer = true,
  fillin/show-answer = true,
  solution/show-solution = show-stay,
}

% ---------- 教师版解答样式 ----------
\newtcolorbox{solutionblock}{
  enhanced, breakable,
  colback=edu-blue-bg,
  colframe=edu-blue!25,
  boxrule=0pt,
  borderline west={2pt}{0pt}{edu-blue!50},
  arc=0pt,
  left=3mm, right=3mm, top=2mm, bottom=2mm,
  before skip=2mm,
  after skip=3mm,
  fontupper=\small,
}

% 解答步骤编号
\newcounter{solstep}
\newcommand{\solstep}[2]{%
  \stepcounter{solstep}%
  \par\medskip
  \noindent\hangindent=6mm
  {\color{edu-blue}\bfseries\textcircled{\scriptsize\thesolstep}\enspace #1}\enspace #2\par
}

\begin{document}

\section*{等腰三角形与直角三角形存在性 · 专题练习（教师版）}

\section*{一、选择题（每题 3 分，共 18 分）}


\begin{question}[points=3]
  在 $x$ 轴上存在点 $P$ 使 $\triangle PAB$ 为等腰三角形。已知 $A(0,3)$，$B(4,0)$，则下列结论正确的是
  \begin{choices}
    \item 只有 $2$ 个符合条件的 $P$ 点
    \item 一定有 $3$ 个符合条件的 $P$ 点
    \item 最多有 $6$ 个符合条件的 $P$ 点（三等腰，每种最多 $2$ 个解）
    \item 条件不足，无法判断
  \end{choices}
  \paren[C]
  \begin{solutionblock}
    等腰三角形有三种等腰情况，每种对应一个一元二次方程（关于 $P$ 的横坐标），一元二次方程最多有两个实数解。因此最多 $3 \times 2 = 6$ 个解。实际解的数量取决于各方程判别式是否大于零，以及是否有退化解需要排除。
  \end{solutionblock}
\end{question}



\begin{question}[points=3]
  已知 $A(-1,2)$，$B(3,2)$，$P$ 在 $y$ 轴上，$\triangle ABP$ 为直角三角形。当 $\angle P = 90°$ 时，$P$ 的纵坐标满足的方程是
  \begin{choices}
    \item $y^2 - 4y + 1 = 0$
    \item $y^2 - 4y + 3 = 0$
    \item $y^2 + 4y - 3 = 0$
    \item $y^2 + 4y + 3 = 0$
  \end{choices}
  \paren[A]
  \begin{solutionblock}
    $P(0,y)$，$\overrightarrow{PA}=(-1,2-y)$，$\overrightarrow{PB}=(3,2-y)$。

$\overrightarrow{PA}\cdot\overrightarrow{PB}=(-1)(3)+(2-y)^2=-3+y^2-4y+4=y^2-4y+1=0$。
  \end{solutionblock}
\end{question}



\begin{question}[points=3]
  $\triangle ABC$ 中 $A(0,0)$，$B(6,0)$，$C(0,8)$。$P$ 在 $x$ 轴上使 $\triangle PAC$ 为等腰三角形。下列哪种情况在 $x$ 轴上一定有解？
  \begin{choices}
    \item $PA = PC$（$P$ 在 $AC$ 的垂直平分线上）
    \item $AP = AC$
    \item $CP = CA$
    \item 三种情况在 $x$ 轴上都有解
  \end{choices}
  \paren[B]
  \begin{solutionblock}
    $PA=PC$：$P$ 在 $AC$ 的垂直平分线上。$AC$ 中点 $(0,4)$，$AC$ 垂直平分线为 $y=4$，与 $x$ 轴平行无交点，无解。

$AP=AC=8$：以 $A(0,0)$ 为圆心、半径 $8$ 的圆与 $x$ 轴交于 $(\pm 8,0)$，有解。

$CP=CA=8$：以 $C(0,8)$ 为圆心、半径 $8$ 的圆方程 $x^2+(y-8)^2=64$，与 $x$ 轴（$y=0$）联立得 $x^2+64=64$，$x=0$，$P(0,0)=A$，退化排除，无解。
  \end{solutionblock}
\end{question}



\begin{question}[points=3]
  直角 $\triangle ABC$ 中 $\angle C = 90°$，$AC=3$，$BC=4$。$P$ 在 $AB$ 上运动，使 $\triangle PBC$ 为等腰三角形。下列哪种情况一定有解？
  \begin{choices}
    \item $PB=PC$（$P$ 在 $BC$ 的垂直平分线上）
    \item $BP=BC=4$
    \item $CP=CB=4$
    \item 以上都不一定
  \end{choices}
  \paren[B]
  \begin{solutionblock}
    $BP=BC=4$：以 $B$ 为圆心、半径 $4$ 画圆。$B$ 在线段 $AB$ 上，$AB=\sqrt{3^2+4^2}=5>4$，圆一定与线段 $AB$ 相交于两点（$B$ 本身和另一侧一点），有解。

$PB=PC$：$P$ 在 $BC$ 的垂直平分线上。$BC$ 中垂线不一定与 $AB$ 相交。

$CP=CB=4$：以 $C$ 为圆心、半径 $4$ 画圆。$C$ 到 $AB$ 的距离为 $\dfrac{AC\cdot BC}{AB}=\dfrac{12}{5}$，圆心到直线距离 $12/5 < 4$，确实相交。但作为选择题，$BP=BC$ 的论证更直接。
  \end{solutionblock}
\end{question}



\begin{question}[points=3]
  已知 $A(2,0)$，$B(0,4)$。$P$ 在 $x$ 轴上使 $\triangle ABP$ 为直角三角形。$\angle A=90°$ 时，$P$ 的坐标为
  \begin{choices}
    \item $P(2,0)$（即 $P=A$，退化）
    \item $P(2,0)$
    \item 不存在符合条件的 $P$
    \item $P$ 有两个解
  \end{choices}
  \paren[C]
  \begin{solutionblock}
    $\angle A=90° \Leftrightarrow \overrightarrow{AB}\cdot\overrightarrow{AP}=0$。

$\overrightarrow{AB}=(-2,4)$，$\overrightarrow{AP}=(x-2,0)$。

$(-2)(x-2)+4\cdot 0=-2(x-2)=0$，$x=2$。

$P(2,0)=A$，三点 $A,B,P$ 退化为线段 $AB$（$P$ 与 $A$ 重合），不构成三角形，排除。

因此 $\angle A=90°$ 时无有效解。
  \end{solutionblock}
\end{question}



\begin{question}[points=3]
  下列关于等腰三角形存在性问题的说法，正确的是
  \begin{choices}
    \item 等腰三角形一定有 $3$ 个解（三等腰各一个）
    \item 求出解后不需要检验三点是否共线
    \item 当动点与某个已知顶点重合时，三角形退化为线段，该解需要排除
    \item 直角三角形只需要讨论两种情况（直角在动点和直角在一个定点）
  \end{choices}
  \paren[C]
  \begin{solutionblock}
    A错：每种等腰情况最多2个解（对应一元二次方程），且可能无解（判别式<0）或退化排除，实际解数不确定。

B错：三点共线时不能构成三角形，必须检验。

C正确：动点与已知顶点重合时三角形退化为线段，该解必须排除。

D错：直角三角形有三种直角情况（三个顶点都可以是直角），需要全部讨论。
  \end{solutionblock}
\end{question}


\section*{二、填空题（每题 3 分，共 18 分）}


\begin{question}[points=3]
  $A(0,0)$，$B(4,0)$，$C(2,3)$。$P$ 在 $x$ 轴上使 $\triangle PBC$ 为等腰三角形。当 $BP=BC$ 时，$P$ 的横坐标为 \fillin。
  \fillin
  \paren[$x=4\pm\sqrt{13}$]
  \begin{solutionblock}
    $BC^2=(2-4)^2+(3-0)^2=4+9=13$，$BC=\sqrt{13}$。

$(x-4)^2+0=13$，$(x-4)^2=13$，$x=4\pm\sqrt{13}$。

验证：$P(4+\sqrt{13},0)$ 和 $P(4-\sqrt{13},0)$ 均不与 $B,C$ 共线。有效。
  \end{solutionblock}
\end{question}



\begin{question}[points=3]
  $A(1,0)$，$B(0,2)$。$P$ 在 $y$ 轴上使 $\triangle ABP$ 为直角三角形。$\angle A=90°$ 时，$P$ 的纵坐标为 \fillin。
  \fillin
  \paren[$y=-\dfrac{1}{2}$]
  \begin{solutionblock}
    $\angle A=90°$：$\overrightarrow{AB}=(-1,2)$，$\overrightarrow{AP}=(-1,y)$。

$\overrightarrow{AB}\cdot\overrightarrow{AP}=(-1)(-1)+2\cdot y=1+2y=0$，$y=-\dfrac{1}{2}$。

$P\left(0,-\dfrac{1}{2}\right)$。验证：$A,B,P$ 不共线（$A(1,0),B(0,2),P(0,-1/2)$，面积不为零）。有效。
  \end{solutionblock}
\end{question}



\begin{question}[points=3]
  正方形 $ABCD$ 中 $A(0,0)$，$B(4,0)$，$C(4,4)$，$D(0,4)$。$P$ 在直线 $y=x$ 上使 $\triangle ABP$ 为等腰三角形。当 $AP=AB$ 时，$P$ 到原点的距离为 \fillin。
  \fillin
  \paren[$4$]
  \begin{solutionblock}
    $P(t,t)$，$AB=4$。

$AP^2=t^2+t^2=2t^2=AB^2=16$，$t^2=8$，$t=\pm 2\sqrt{2}$。

$OP=\sqrt{t^2+t^2}=\sqrt{16}=4$。

验证：$P(2\sqrt{2},2\sqrt{2})$ 和 $P(-2\sqrt{2},-2\sqrt{2})$ 均不与 $A,B$ 共线。有效。
  \end{solutionblock}
\end{question}



\begin{question}[points=3]
  已知 $A(-2,0)$，$B(2,0)$，$P$ 在直线 $y=x+1$ 上使 $\triangle ABP$ 为直角三角形。$\angle P=90°$ 时，$P$ 的坐标满足的方程为 $\underline{\quad\quad}$（化为一般式）。
  \fillin
  \paren[$2x^2+2x-3=0$]
  \begin{solutionblock}
    $\angle P=90°$：$P$ 在以 $AB$ 为直径的圆上。$AB$ 中点 $(0,0)$，半径 $2$。

圆方程：$x^2+y^2=4$。

联立 $y=x+1$：$x^2+(x+1)^2=4$，$2x^2+2x-3=0$。

$\Delta=4+24=28>0$，有两个实数解。
  \end{solutionblock}
\end{question}



\begin{question}[points=3]
  矩形 $ABCD$ 中 $A(0,0)$，$B(6,0)$，$C(6,4)$，$D(0,4)$。$P$ 在 $BC$ 上使 $\triangle APD$ 为等腰三角形。当 $AP=DP$ 时，$P$ 的坐标为 \fillin。
  \fillin
  \paren[$P(6,2)$]
  \begin{solutionblock}
    $AP=DP$：$P$ 在 $AD$ 的垂直平分线上。$AD$ 中点 $(0,2)$，$AD$ 垂直平分线为 $y=2$。

$P$ 在 $BC$ 上（$x=6$），$y=2$ 与 $x=6$ 的交点 $P(6,2)$。

验证：$AP^2=36+4=40=DP^2$。$\checkmark$
  \end{solutionblock}
\end{question}



\begin{question}[points=3]
  $\triangle ABC$ 中 $A(0,0)$，$B(6,0)$，$C(3,4)$。$P$ 在 $y$ 轴上使 $\triangle PBC$ 为等腰三角形。$PC=BC$ 时，$P$ 的坐标为 \fillin。
  \fillin
  \paren[$P(0,0)$ 或 $P(0,8)$]
  \begin{solutionblock}
    $BC^2=(3-6)^2+(4-0)^2=9+16=25$，$BC=5$。

$P(0,y)$，$PC^2=(0-3)^2+(y-4)^2=9+(y-4)^2$。

$PC=BC$：$9+(y-4)^2=25$，$(y-4)^2=16$，$y-4=\pm 4$，$y=0$ 或 $y=8$。

$P(0,0)=A$：$A(0,0),B(6,0),C(3,4)$ 构成三角形（面积 $=12\neq 0$），$\triangle PBC=\triangle ABC$ 有效。
$P(0,8)$：不共线，有效。
  \end{solutionblock}
\end{question}


\section*{三、解答题（每题 10 分，共 30 分）}

\needspace{8\baselineskip}

\begin{problem}[points=10]
  直线 $y=-\dfrac{4}{3}x+4$ 与 $x$ 轴、$y$ 轴分别交于 $A$、$B$ 两点。$P$ 是 $x$ 轴上的动点，$\triangle ABP$ 为等腰三角形。

（1）求 $A$、$B$ 的坐标和 $AB$ 的长；

（2）求所有符合条件的 $P$ 的坐标。
  \answerarea[80mm]
  \setcounter{solstep}{0}
  \begin{solutionblock}
    \textbf{答案：}$P\left(\dfrac{3}{2},0\right)$，$P(8,0)$，$P(-2,0)$，$P(-3,0)$\par\medskip
    \solstep{确定 $A,B$ 坐标和 $AB$ 长}{$y=-\dfrac{4}{3}x+4$ 与 $x$ 轴交点 $A(3,0)$，与 $y$ 轴交点 $B(0,4)$。$AB=5$。}
    \solstep{设 $P(x,0)$，分类讨论}{三等腰：$AP=BP$、$AP=AB=5$、$BP=AB=5$。}
    \solstep{情形一 $AP=BP$}{$(x-3)^2+16=x^2+16$，$x=\dfrac{3}{2}$。}
    \solstep{情形二 $AP=AB=5$}{$(x-3)^2=25$，$x=8$ 或 $x=-2$。}
    \solstep{情形三 $BP=AB=5$}{$x^2+16=25$，$x=3$（退化排除）或 $x=-3$。}
    \solstep{汇总}{$P\left(\dfrac{3}{2},0\right)$，$P(8,0)$，$P(-2,0)$，$P(-3,0)$。}
  \end{solutionblock}
\end{problem}


\needspace{8\baselineskip}

\begin{problem}[points=10]
  正方形 $OABC$ 边长 $4$，$O$ 在原点，$A$ 在 $x$ 轴正方向，$C$ 在 $y$ 轴正方向。$D(2,0)$。$P$ 从 $O$ 沿 $O \to A \to B$ 运动（$1$ 单位/秒）。

是否存在 $P$ 使 $\triangle CDP$ 为等腰三角形？若存在，求所有 $P$ 的坐标。
  \answerarea[80mm]
  \setcounter{solstep}{0}
  \begin{solutionblock}
    \textbf{答案：}参考例题12的解题结构。$CD=2\sqrt{5}$。分 $O\to A$ 段和 $A\to B$ 段参数化 $P$，三等腰分类讨论。\par\medskip
    $O(0,0)$，$A(4,0)$，$B(4,4)$，$C(0,4)$，$D(2,0)$。$CD^2=4+16=20$，$CD=2\sqrt{5}$。

$O\to A$ 段：$P(t,0)$，$0\leq t\leq 4$。
$A\to B$ 段：$P(4,t-4)$，$4\leq t\leq 8$。

三等腰 $CD=DP$、$CP=DP$、$CD=CP$ 在每段上分别求解。

$O\to A$ 段，$CD=DP$：$(t-2)^2+20=20$，$(t-2)^2=0$，$t=2$，$P(2,0)=D$，退化排除。

$O\to A$ 段，$CP=DP$：$t^2+16=(t-2)^2$，$t^2+16=t^2-4t+4$，$4t=-12$，$t=-3<0$，越界。

$O\to A$ 段，$CD=CP$：$t^2+16=20$，$t^2=4$，$t=\pm 2$。$t=2$：$P(2,0)=D$，退化排除。$t=-2<0$，越界。

$A\to B$ 段：$P(4,t-4)$，$4\leq t\leq 8$。

$CD=DP$：$(4-2)^2+(t-4)^2=20$，$4+(t-4)^2=20$，$(t-4)^2=16$，$t-4=\pm 4$，$t=0$（越界）或 $t=8$。$P(4,4)=B$。

$CP=DP$：$16+(t-4-4)^2=4+(t-4)^2$，$16+(t-8)^2=4+(t-4)^2$，$16+t^2-16t+64=4+t^2-8t+16$，$-16t+80=-8t+20$，$-8t=-60$，$t=7.5 \in [4,8]$。$P(4,3.5)$。

$CD=CP$：$16+(t-8)^2=20$，$(t-8)^2=4$，$t-8=\pm 2$，$t=10$（越界）或 $t=6$。$P(4,2)$。

汇总：$P(4,4)$，$P(4,3.5)$，$P(4,2)$。共 $3$ 个解。
  \end{solutionblock}
\end{problem}


\needspace{8\baselineskip}

\begin{problem}[points=10]
  $\triangle ABC$ 中 $AB=AC=5$，$BC=6$。$D$ 是 $AB$ 上一点，$DE \parallel BC$ 交 $AC$ 于 $E$。以 $DE$ 为边在 $DE$ 下方作正方形 $DEFG$（$G$ 在 $\triangle ABC$ 内部）。

当 $\triangle BDG$ 为等腰三角形时，求 $AD$ 的长。
  \answerarea[80mm]
  \setcounter{solstep}{0}
  \begin{solutionblock}
    \textbf{答案：}$AD = \dfrac{20}{7}$ 或 $AD = \dfrac{25}{11}$ 或 $AD = \dfrac{125}{73}$\par\medskip
    \solstep{建系}{$B(-3,0)$，$C(3,0)$，$A(0,4)$。}
    \solstep{参数化 $D$、$E$}{$AD=t$。$D\left(-\dfrac{3t}{5}, 4-\dfrac{4t}{5}\right)$，$E\left(\dfrac{3t}{5}, 4-\dfrac{4t}{5}\right)$。$DE=\dfrac{6t}{5}$。}
    \solstep{确定 $G$ 的坐标}{$G\left(-\dfrac{3t}{5}, 4-2t\right)$。}
    \solstep{计算 $BG^2$}{$BG^2 = \dfrac{109t^2-490t+625}{25}$。}
    \solstep{三等腰分类讨论}{$BD=BG$：$t=\dfrac{20}{7}$；$BD=DG$：$t=\dfrac{25}{11}$；$BG=DG$：$t=\dfrac{125}{73}$。}
    \solstep{汇总}{$AD = \dfrac{20}{7}$ 或 $\dfrac{25}{11}$ 或 $\dfrac{125}{73}$。}
  \end{solutionblock}
\end{problem}


\clearpage
\section*{参考答案}


\begin{keyideabox}
  \textbf{选择题}
  1. C

2. A

3. B

4. B

5. C

6. C
\end{keyideabox}



\begin{keyideabox}
  \textbf{填空题}
  1. $x=4\pm\sqrt{13}$

2. $y=-\dfrac{1}{2}$

3. $4$

4. $2x^2+2x-3=0$

5. $P(6,2)$

6. $P(0,0)$ 或 $P(0,8)$
\end{keyideabox}



\begin{keyideabox}
  \textbf{解答题}
  第 1 题：$P\left(\dfrac{3}{2},0\right)$，$P(8,0)$，$P(-2,0)$，$P(-3,0)$，共 $4$ 个解。

第 2 题：$P(4,4)$，$P(4,3.5)$，$P(4,2)$，共 $3$ 个解。

第 3 题：$AD = \dfrac{20}{7}$ 或 $AD = \dfrac{25}{11}$ 或 $AD = \dfrac{125}{73}$。
\end{keyideabox}



\end{document}
